\documentclass{article}
\usepackage{amsmath,amssymb,amsthm,mathtools}
\newtheorem{theorem}{Theorem}
\newtheorem{lemma}{Lemma}
\newtheorem{proposition}{Proposition}
\newtheorem{corollary}{Corollary}
\newtheorem{definition}{Definition}
\newtheorem{example}{Example}
\title{ShadowBench probability/L3/prob\_gen\_L3\_001}
\begin{document}
\maketitle
% Problem id: probability/L3/prob_gen_L3_001
% Companion instructions: docs/instructions.md
% Lean target: ShadowBench/Source/Main.lean

Theorem(`BuffonNeedle`).

Suppose a short needle of length $\ell$ is dropped onto a paper ruled with equally spaced lines of distance $d \ge \ell$. Then the probability that the needle crosses a line on the paper is $\frac{2\ell}{\pi d}$.

Proof.

We can assume the uniform probability distribution over $[0, d] \times [0, \pi]$, where each component means the midpoint position relative to the closest lower line and the angle relative to a fixed direction of the ruling respectively. Then the needle crosses a line with angle $\theta$ exactly when its height $y$ of midpoint is lower than $\frac{\ell \sin \theta}{2}$ or higher than $d - \frac{\ell \sin \theta}{2}$. Thus the marginal probability for $\theta$ is $\frac{\ell \sin \theta}{d}$. Now we get the probability as the average of probability over $\theta$, namely
$$
\frac{1}{\pi} \int_0^\pi \frac{\ell \sin \theta}{d} d\theta = \frac{2\ell}{\pi d}.
$$
\end{document}
